\documentclass[aspectratio=169,10pt]{beamer}
% MyPPTSkill managed-font-entry
\usepackage{iftex}
\ifPDFTeX
  \PackageError{beamer-starter}{Use XeLaTeX or LuaLaTeX}{The font setup uses OpenType files.}
\fi
\usepackage{fontspec}
\usepackage[fontset=none]{ctex}
\InputIfFileExists{myppt-fonts.tex}{}{%
  \PackageError{beamer-starter}{Managed font configuration is missing}{Build with MyPPTSkill build_beamer.py or run setup_fonts.py for this project first.}}

\usepackage{amsmath,booktabs,tikz,pgfplots}
\usepgfplotslibrary{groupplots}
\pgfplotsset{compat=1.18}
\definecolor{Ink}{HTML}{17324D}
\definecolor{Accent}{HTML}{167C80}
\definecolor{Compare}{HTML}{BC5D28}
\definecolor{Muted}{HTML}{566778}
\definecolor{Pale}{HTML}{EDF4F6}
\setbeamercolor{normal text}{fg=Ink,bg=white}
\setbeamercolor{frametitle}{fg=Ink,bg=white}
\setbeamercolor{structure}{fg=Accent}
\setbeamercolor{block title}{fg=Ink,bg=Pale}
\setbeamercolor{block body}{fg=Ink,bg=Pale}
\setbeamercolor{footline}{fg=Muted}
\setbeamertemplate{navigation symbols}{}
\setbeamertemplate{blocks}[default]
\setbeamertemplate{itemize item}{\small\raisebox{0.1ex}{\textbullet}}
\setbeamersize{text margin left=7mm,text margin right=7mm}

% Semantic font roles; TeX pt on the native 16 x 9 cm page.
\newcommand{\BodyFont}{\fontsize{10.5}{13.4}\selectfont}
\newcommand{\AuxFont}{\fontsize{9.2}{11.7}\selectfont}
\newcommand{\AuxNote}[1]{{\AuxFont\color{black!62}#1\par}}
\newcommand{\PlotFont}{\fontsize{8.8}{10.5}\selectfont}
\newcommand{\CodeFont}{\fontsize{9.2}{11.8}\selectfont}
\setbeamerfont{normal text}{size=\BodyFont}
\setbeamerfont{frametitle}{size=\fontsize{14.5}{18}\selectfont,series=\mdseries}
\setbeamerfont{block title}{size=\BodyFont,series=\mdseries}
\setbeamerfont{block body}{size=\BodyFont}
\setbeamerfont{footline}{size=\fontsize{8}{10}\selectfont}
\setbeamertemplate{frametitle}{%
  \vspace{2mm}\usebeamerfont{frametitle}\insertframetitle\par\vspace{1.5mm}}
% Count physical PDF pages, including every overlay and appendix page.
\setbeamertemplate{footline}{%
  \hfill\usebeamerfont{footline}\usebeamercolor[fg]{footline}%
  \insertpagenumber/\insertdocumentendpage\hspace{7mm}\vspace{2mm}}
\AtBeginDocument{\BodyFont
  \setlength{\abovedisplayskip}{4pt}\setlength{\belowdisplayskip}{4pt}
  \setlength{\abovedisplayshortskip}{3pt}\setlength{\belowdisplayshortskip}{3pt}}
\newcommand{\Key}[1]{{\color{Accent}\mdseries #1}}
\newcommand{\Insight}[1]{\par\vspace{2mm}%
  \begin{beamercolorbox}[sep=1.8mm,wd=\textwidth]{block body}\BodyFont #1\end{beamercolorbox}}

\pgfplotsset{starter axis/.style={
  width=6.8cm,height=4.65cm,
  axis line style={Ink!55},tick style={Ink!55},
  grid=major,grid style={Ink!10},
  tick label style={font=\PlotFont},
  label style={font=\PlotFont},title style={font=\PlotFont},
  legend style={font=\PlotFont,draw=none,fill=white,inner sep=2pt},
  line width=0.8pt,mark size=1.6pt,
  scaled ticks=false,
}}

% Load after fonts.tex and theme.tex. Values are starting points for visual review.
\usetheme{Boadilla}
\usecolortheme{beaver}
\setbeamertemplate{navigation symbols}{}
\setbeamersize{text margin left=.6cm,text margin right=.6cm}
% fonts.tex loads the exact Fandol/Latin Modern files from the managed bundle.
\definecolor{Ink}{HTML}{282B30}
\definecolor{Accent}{HTML}{93232C}
\definecolor{Pale}{HTML}{F5F5F5}
\setbeamercolor{normal text}{fg=Ink,bg=white}
\setbeamercolor{frametitle}{fg=Accent,bg=gray!9}
\setbeamercolor{title}{fg=Accent,bg=gray!14}
\setbeamercolor{structure}{fg=Accent}
\setbeamercolor{block title}{fg=Accent,bg=gray!14}
\setbeamercolor{block body}{fg=Ink,bg=Pale}
\setbeamerfont{normal text}{size=\BodyFont,series=\mdseries}
% Calibrate long titles in the actual cover; these are initial sizes.
\setbeamerfont{title}{size=\fontsize{19}{24}\selectfont,series=\mdseries}
\setbeamerfont{subtitle}{size=\fontsize{12}{16}\selectfont,series=\mdseries}
\setbeamerfont{frametitle}{size=\fontsize{14.5}{18}\selectfont,series=\mdseries}
\setbeamerfont{block title}{size=\BodyFont,series=\mdseries}
\setbeamertemplate{frametitle}{%
  \nointerlineskip
  % Single-line 14.5pt FandolHei calibration. Recheck visible ink after title/font changes.
  \begin{beamercolorbox}[wd=\paperwidth,sep=0pt,leftskip=.6cm,rightskip=.6cm,ht=.57cm,dp=.24cm]{frametitle}
    \usebeamerfont{frametitle}\insertframetitle
  \end{beamercolorbox}}
\setbeamercolor{author in head/foot}{fg=white,bg=Accent}
\setbeamercolor{title in head/foot}{fg=Accent,bg=gray!14}
\setbeamercolor{date in head/foot}{fg=Accent,bg=gray!23}
% Set this only when an author is supplied. Otherwise the title occupies its space.
\newcommand{\BioFooterAuthor}{}
\newcommand{\BioFooterCell}[3]{%
  \begingroup
  % Keep all background edges aligned; optical correction moves text only.
  \begin{beamercolorbox}[wd=#1,ht=.26cm,dp=.10cm,center]{#2}
    \ifstrequal{#2}{date in head/foot}{%
      \raisebox{.28pt}[0pt][0pt]{\fontsize{7.9}{9}\selectfont\mdseries #3}%
    }{\fontsize{7.9}{9}\selectfont\mdseries #3}%
  \end{beamercolorbox}\endgroup}
\setbeamertemplate{footline}{%
  \leavevmode\hbox{%
    \ifdefempty{\BioFooterAuthor}{%
      \BioFooterCell{.78\paperwidth}{title in head/foot}{\insertshorttitle}%
    }{%
      \BioFooterCell{.28\paperwidth}{author in head/foot}{\BioFooterAuthor}%
      \BioFooterCell{.50\paperwidth}{title in head/foot}{\insertshorttitle}%
    }%
    \BioFooterCell{.22\paperwidth}{date in head/foot}{\insertpagenumber/\insertdocumentendpage}%
  }\vskip0pt}

\title{二次函数的步长与收敛}
\hypersetup{pdftitle={二次函数的步长与收敛},pdfauthor={},pdfsubject={梯度下降的可复算示例},pdfcreator={},pdfproducer={}}

\begin{document}

\begin{frame}[c,label=recurrence]{步长决定误差每轮保留多少}
\begin{columns}[T,onlytextwidth]
\begin{column}{0.49\textwidth}
\Key{问题：从 $x_0=4$ 找到最小值}
\[
 f(x)=\tfrac12 x^2,\qquad x^*=0
\]
梯度是 $f'(x)=x$。每轮沿负梯度更新，步长 $\eta$ 控制移动距离：
\begin{align*}
 x_{k+1}&=x_k-\eta f'(x_k)\\
        &=(1-\eta)x_k,\\
 x_k&=4(1-\eta)^k.
\end{align*}
当 $0<\eta<2$ 时，$|1-\eta|<1$，误差逐轮缩小。
\end{column}
\begin{column}{0.46\textwidth}
\Key{首两轮：同一起点，两种收缩速度}
\medskip

\renewcommand{\arraystretch}{1.25}
\begin{tabular}{@{}lrr@{}}
\toprule
迭代 $k$ & $\eta=0.25$ & $\eta=0.75$\\
\midrule
0 & 4 & 4\\
1 & 3 & 1\\
2 & 2.25 & 0.25\\
\bottomrule
\end{tabular}
\medskip

表内为 $x_k$；每轮分别保留上一轮误差的 $75\%$ 与 $25\%$。
\medskip

目标值随误差的平方缩小：
\[
 f(x_{k+1})=(1-\eta)^2 f(x_k).
\]
\end{column}
\end{columns}
\Insight{在这个曲率为 1 的二次函数上，$\eta=0.75$ 比 $\eta=0.25$ 更快逼近最优点。}
\end{frame}

\begin{frame}[c,label=curves]{距离与目标值呈现不同的收敛尺度}
\centering
\begin{tikzpicture}
\begin{groupplot}[
  starter axis,
  group style={group size=2 by 1,horizontal sep=17mm},
  xmin=0,xmax=6,xtick={0,2,4,6},xlabel={迭代次数 $k$},
  samples at={0,1,2,3,4,5,6},
]
\nextgroupplot[title={(a) 到最优点的距离},ylabel={$|x_k|$},
  ymin=0,ymax=4.3,ytick={0,1,2,3,4},legend pos=north east]
\addplot[Compare,mark=square*] {4*pow(0.75,x)};
\addlegendentry{$\eta=0.25$}
\addplot[Accent,mark=*] {4*pow(0.25,x)};
\addlegendentry{$\eta=0.75$}
\nextgroupplot[title={(b) 目标值的数量级},ylabel={$f(x_k)$},
  ymode=log,log basis y=10,ymin=1e-7,ymax=20,
  ytick={1e-6,1e-4,1e-2,1},
  yticklabels={$10^{-6}$,$10^{-4}$,$10^{-2}$,$10^0$}]
\addplot[Compare,mark=square*] {8*pow(0.75,2*x)};
\addplot[Accent,mark=*] {8*pow(0.25,2*x)};
\end{groupplot}
\end{tikzpicture}

\par\vspace{1mm}
\raggedright
\AuxNote{图中点由 $x_k=4(1-\eta)^k$、$f(x_k)=8(1-\eta)^{2k}$ 计算。}
\par\medskip
\BodyFont
左图看绝对距离；右图用对数轴显示数量级，相邻主刻度对应 $100$ 倍的变化。
\Insight{$k=4$ 时，两种步长的距离分别为 $1.265625$ 和 $0.015625$；对应目标值相差 $6561$ 倍。}
\end{frame}

\begin{frame}[c,fragile,label=stopping]{同一精度门槛对应 13 次与 3 次更新}
\begin{columns}[T,onlytextwidth]
\begin{column}{0.49\textwidth}
\Key{停止条件：$|x_k|\leq 0.1$}
\medskip

对 $0<\eta<1$，由递推式直接得到：
\[
 k_{\min}=\left\lceil\frac{\log(0.1/4)}{\log(1-\eta)}\right\rceil.
\]
\renewcommand{\arraystretch}{1.2}
\begin{tabular}{@{}l@{\hspace{4mm}}r@{\hspace{4mm}}r@{}}
\toprule
步长 & 更新次数 & 最终 $|x_k|$\\
\midrule
$0.25$ & $13$ & $0.0950$\\
$0.75$ & $3$ & $0.0625$\\
\bottomrule
\end{tabular}
\medskip

每轮计算一次梯度时，梯度调用量从 13 次降至 3 次，减少约 $76.9\%$。
\end{column}
\begin{column}{0.46\textwidth}
\Key{用短循环复算停止时刻}
\medskip

{\CodeFont
\begin{verbatim}
for eta in (0.25, 0.75):
    x, k = 4.0, 0
    while abs(x) > 0.1:
        x = (1 - eta) * x
        k += 1
    print(eta, k, abs(x))
\end{verbatim}
}
\medskip
$\eta=0.25$ 的第 12 轮仍有 $|x_{12}|\approx0.1267$；$\eta=0.75$ 的第 2 轮仍为 $0.25$，均尚未达到门槛。
\end{column}
\end{columns}
\Insight{步长选择取决于曲率：若 $f(x)=\tfrac12 a x^2$ 且 $a>0$，收敛范围变为 $0<\eta<2/a$。}
\end{frame}

\end{document}

